\section{Introduction}

MedGuard-CXR is a safety-aware chest X-ray engineering prototype. The project combines a DenseNet121 image classifier, calibration and abstention logic, pneumonia-specific RSNA grounding checks, and a rule-based VQA/API/demo layer that keeps research-only disclaimers visible.

The purpose of this report is to publish the current real NIH classifier results and to preserve the continuation plan for calibration, RSNA grounding, and safe VQA work. The system is not a clinical diagnostic tool, has not been prospectively evaluated, and must not be used for patient care.

The current repository status is explicitly \textbf{REAL COMPLETED}: NIH classifier training, NIH evaluation, calibration, and RSNA Pneumonia grounding artifacts exist. By contrast, the \textbf{IDEALIZED FUTURE / NOT YET RUN} path includes broader VLM/QLoRA training, free-text VLM validation, clinical validation, and any other aspirational model-training milestones that are not present in the current repository artifacts. Phase 4A VQA/API/Gradio remains rule-based and smoke-tested. Phase 4B VLM/QLoRA is deferred and not trained.

\subsection{Research Questions}

\textbf{RQ1: does per-class post-hoc calibration meaningfully reduce expected calibration error (ECE) under multi-label class imbalance with noisy, silver-standard labels?} Answer: no. Temperature scaling moves macro ECE from 0.3144 to 0.3113 (\texttt{results/calibration\_report.json}), an improvement small enough to report as a negative result rather than a success. This is reported as a structural finding, not a tuning failure: temperature scaling is a single positive scalar divisor and is rank-preserving, so it can sharpen or flatten a probability distribution but cannot correct a label-ordering error, and NIH ChestX-ray14's silver-standard labels are themselves NLP-mined and noisy, which caps what any monotonic recalibration can fix.

\textbf{RQ2: how far do point-based and IoU-based grounding metrics diverge under cross-dataset transfer (NIH $\rightarrow$ RSNA)?} Answer: substantially, and the divergence separates cleanly into a fixable and an unfixable component. Sweeping the Grad-CAM extraction threshold on the RSNA validation split (Section~\ref{sec:results}) leaves pointing-game accuracy unchanged at 0.5138 across all seven thresholds tested, while mAP@0.5 improves $10.8\times$ over the same range. Point-based and IoU-based metrics are therefore not two views of the same error: the former tracks whether the CAM peak is in the right place, the latter also demands the right box shape, and only the shape component responds to post-hoc threshold tuning.

\textbf{RQ3: does a classifier-consistency gate suppress hallucination in a generative presentation layer?} Answer: not yet established either way. The rule-based VQA baseline reports a hallucination rate of 0.0 (\texttt{results/vlm\_zero\_shot\_eval.json}), but this baseline is deterministic template generation over a synthetic, template-generated QA set, not a generative model, so a hallucination rate of zero is close to true by construction rather than evidence that the gate suppresses anything. The free-text VLM zero-shot backend, where a consistency gate would actually have generative behavior to constrain, has not been executed in this environment (\texttt{bitsandbytes} unavailable), so the question this RQ actually asks remains open.
